\documentclass[9pt,a4paper]{article}

\usepackage[utf8]{inputenc}
\usepackage[T1]{fontenc}
\usepackage[polish]{babel}
\usepackage{lmodern}

\usepackage[top=0.55cm,bottom=0.55cm,left=0.65cm,right=0.65cm]{geometry}
\usepackage{xcolor}
\usepackage{microtype}
\usepackage{enumitem}
\usepackage{tabularx}
\usepackage{array}
\usepackage[hidelinks]{hyperref}

\pagestyle{empty}
\setlength{\parindent}{0pt}
\setlength{\parskip}{0pt}
\setlength{\tabcolsep}{0pt}
\renewcommand{\familydefault}{\sfdefault}
\emergencystretch=2em
\raggedbottom

\definecolor{navy}{HTML}{172A46}
\definecolor{teal}{HTML}{167C80}
\definecolor{gold}{HTML}{D39A3A}
\definecolor{burgundy}{HTML}{7A304A}
\definecolor{pale}{HTML}{F1F5F5}
\definecolor{ink}{HTML}{27313D}
\definecolor{muted}{HTML}{65717D}

\newcommand{\cvsection}[2]{%
  \vspace{0.40em}%
  {\color{#1}\bfseries\fontsize{9.4}{10.6}\selectfont
    \MakeUppercase{#2}\par}%
  \vspace{0.06em}%
  {\color{#1}\hrule height 0.7pt}%
  \vspace{0.14em}%
}

\newcommand{\sideitem}[2]{%
  {\color{navy}\bfseries #1}\par
  {\color{ink}#2\par}
  \vspace{0.13em}%
}

\newcommand{\cvbullet}[1]{%
  \noindent{\color{gold}\textbullet}\hspace{0.42em}#1\par
  \vspace{0.13em}%
}

\newcommand{\cventry}[4]{%
  \begin{tabularx}{\linewidth}{@{}>{\raggedright\arraybackslash}X
                                  >{\raggedleft\arraybackslash}p{2.15cm}@{}}
    {\fontsize{9}{10.2}\selectfont\bfseries\color{navy}#2} &
    {\fontsize{8}{9.2}\selectfont\bfseries\color{teal}#1}
  \end{tabularx}
  {\fontsize{8.2}{9.4}\selectfont\itshape\color{burgundy}#3\par}
  \vspace{0.06em}
  {\fontsize{8.45}{9.9}\selectfont #4\par}
  \vspace{0.17em}
}

\begin{document}
\color{ink}

% Nagłówek
\noindent
\colorbox{navy}{%
  \parbox{\dimexpr\textwidth-2\fboxsep\relax}{%
    \vspace{0.34em}
    {\fontsize{24}{26}\selectfont\bfseries\color{white} Józef Szwejk}
    \hfill
    {\fontsize{8.5}{10}\selectfont\bfseries\color{gold}
      PROFIL GOŚCINNEGO PRELEGENTA}\par
    \vspace{0.10em}
    {\fontsize{9.2}{11}\selectfont\color{white}
      Doświadczenia zwykłego człowieka jako żywy punkt wyjścia
      do rozmowy o historii, społeczeństwie i literaturze}\par
    \vspace{0.18em}
    {\color{gold}\rule{\linewidth}{1.1pt}}\par
    \vspace{0.20em}
  }%
}

\vspace{0.26em}

% Lewa kolumna
\noindent
\begin{minipage}[t]{0.295\textwidth}
\vspace{0pt}

\colorbox{pale}{%
  \begin{minipage}[t]{\dimexpr\linewidth-2\fboxsep\relax}
  \vspace{0.14em}
  \raggedright
  \fontsize{8.45}{9.9}\selectfont

  \cvsection{teal}{W skrócie}
  Praktyk, obserwator i gawędziarz. Dzięki doświadczeniu w handlu,
  podróżach i służbie wojskowej potrafię opowiedzieć o wielkich
  wydarzeniach z perspektywy ich codziennego wpływu na człowieka,
  jego języka i zachowania.

  Nie jestem badaczem akademickim, lecz kimś, kto doświadczał
  instytucji od środka i potrafi przetłumaczyć abstrakcyjne pojęcia
  na konkretne sytuacje, pytania i anegdoty.

  \cvsection{teal}{Obszary tematyczne}
  \cvbullet{codzienność monarchii austro-węgierskiej}
  \cvbullet{wojsko, hierarchia i mechanizmy posłuszeństwa}
  \cvbullet{biurokracja widziana od strony obywatela}
  \cvbullet{humor jako sposób komentowania rzeczywistości}
  \cvbullet{relacje jednostki z państwem i instytucją}
  \cvbullet{język rozkazów i język potocznej rozmowy}
  \cvbullet{praca, podróże i mobilność przed I wojną światową}

  \cvsection{teal}{Atuty dla wykładowcy}
  \cvbullet{barwne przykłady zamiast oderwanej abstrakcji}
  \cvbullet{naturalne łączenie historii z literaturą i kulturą}
  \cvbullet{zdolność interpretacji ludzkich reakcji i motywacji}
  \cvbullet{spokój, elastyczność i dystans w trudnych sytuacjach}
  \cvbullet{aktywizowanie słuchaczy przez pytania i anegdoty}
  \cvbullet{kontakt z osobami o różnych doświadczeniach i statusach}

  \cvsection{teal}{Języki}
  \sideitem{Czeski}{Ojczysty, naturalny i potoczny}
  \sideitem{Niemiecki}{Komunikatywny, również w służbie}
  \sideitem{Węgierski}{Znajomość podstawowa}

  \cvsection{teal}{Dane i kontakt}
  \sideitem{Józef Szwejk}{Profil osoby oraz doświadczenie praktyczne}
  Narodowość: czeska\\
  Praga, ul. Bojiště 14\\
  \href{mailto:jozef.szwejk@example.com}
    {jozef.szwejk@example.com}

  \cvsection{teal}{Poza służbą}
  Kynologia, literatura, obserwowanie ludzkich charakterów,
  marsze, życie gospód oraz opowiadanie zdarzeń z różnych epizodów życia.

  \vspace{0.14em}
  \end{minipage}%
}

\end{minipage}
\hfill
% Prawa kolumna
\begin{minipage}[t]{0.685\textwidth}
\vspace{0pt}
\fontsize{8.55}{10}\selectfont

\cvsection{navy}{Propozycja wykładu gościnnego}

{\color{navy}\bfseries\fontsize{10.2}{11.5}\selectfont
Wielka historia z perspektywy zwykłego człowieka}\par
\vspace{0.07em}

Jak opowiadać o epoce przez codzienne doświadczenia: rozkazy,
podróże, rozmowy, pomyłki i spotkania z urzędami. Anegdota może
stać się punktem wyjścia do dyskusji o władzy, języku, hierarchii,
społecznym pochodzeniu i sposobach radzenia sobie z rzeczywistością.

Spotkanie może łączyć krótką narrację z analizą sytuacji, interpretacją
reakcji bohaterów i pytaniem o to, jaką wiedzę o epoce daje nam życie
jednostki, a nie tylko dokument urzędowy czy oficjalne sprawozdanie.

\cvsection{burgundy}{Doświadczenie zawodowe i społeczne}

\cventry
  {1914--1918}
  {Służba w 91. Pułku Piechoty}
  {Ordynans; bezpośredni kontakt z życiem wojskowym i administracją}
  {Wykonywanie poleceń, przekazywanie meldunków i współpraca
  z przełożonymi oraz żołnierzami o różnych pochodzeniach i doświadczeniach.
  Widziałem, jak formalne rozkazy przechodzą przez codzienne relacje,
  jak instytucja działa przez ludzi i jak techniczne polecenie
  staje się źródłem nieporozumień, humoru lub nowych problemów.
  Materiał ten pozwala omawiać dyscyplinę, strukturę, komunikację
  oraz konflikt między procedurą a ludzką interpretacją.}

\cventry
  {1910--1914}
  {Kynologiczny Instytut we współczesnej Pradze}
  {Właściciel; handel, negocjacje i bezpośrednia obsługa klientów}
  {Prowadzenie rozmów handlowych, rozpoznawanie oczekiwań klientów
  i dobieranie sposobu prezentacji do konkretnego rozmówcy.
  Doświadczenie to pozwala mówić o dawnym mieście jako miejscu pracy,
  konsumpcji i codziennej komunikacji. Pokazuje także znaczenie
  wyglądu, tonu głosu, uprzejmości i umiejętności obserwacji
  w uzyskaniu zaufania.}

\cvsection{burgundy}{Epizody jako materiał do analizy}

\cventry
  {1915}
  {Marsz przez południowe Czechy}
  {Samodzielne działanie w zmiennych warunkach}
  {Poszukiwanie macierzystego pułku, orientacja w terenie i kontakt
  z mieszkańcami. Epizod odsłania różnicę między formalnym celem
  instytucji a doświadczeniem osoby próbującej go osiągnąć.
  Można go wykorzystać do dyskusji o przestrzeni, mobilności,
  informacjach niepełnych i sposobach podejmowania decyzji.}

\cventry
  {1914}
  {Droga do garnizonu}
  {Realizacja zadania w warunkach nieporozumienia i presji}
  {Wydarzenia związane z Sarajewem wciągnęły zwykłego mieszkańca
  Pragi w narrację wielkich wydarzeń. To przykład, jak przypadek,
  interpretacja i nagłe podejście urzędowe mogą zmienić los jednostki.
  Pomaga rozmowie o relacji między wojną a codziennym życiem.}

\cvsection{teal}{Przebieg możliwego spotkania}

\begin{tabularx}{\linewidth}{@{}>{\color{gold}\bfseries}p{0.75cm}
                                  >{\raggedright\arraybackslash}X@{}}
  1. & Krótka opowieść i osadzenie epizodu w realiach epoki\\
  2. & Interpretacja zachowań, języka i struktur społecznych\\
  3. & Porównanie doświadczenia jednostki z oficjalną narracją\\
  4. & Zadania grupowe: pytania, alternatywny przebieg zdarzeń\\
  5. & Dyskusja o człowieku, instytucji i odpowiedzialności
\end{tabularx}

\cvsection{navy}{Potencjalne tematy dla zajęć}

{\fontsize{8.2}{9.6}\selectfont
\textbf{Historia codzienności.}
Jak wyglądają wielkie przemiany, gdy obserwuje się je z poziomu
konkretnego człowieka, ulicy, gospody i garnizonu.

\vspace{0.07em}
\textbf{Literatura i świadomość społeczna.}
Jak anegdota, humor i pozorna prostota narracji mogą ujawnić mechanizmy
władzy, kulturę biurokracji i sposób postrzegania świata.

\vspace{0.07em}
\textbf{Władza i posłuszeństwo.}
Jak rozkaz staje się praktyką społeczną oraz dlaczego jego wykonanie
nie zawsze oznacza takie samo rozumienie polecenia.

\vspace{0.07em}
\textbf{Pamięć i interpretacja.}
W jaki sposób doświadczenie jednostki uzupełnia źródła urzędowe
i oficjalne relacje historyczne.
}

\cvsection{navy}{Przygotowanie i ograniczenia}

Praktyczne doświadczenie handlowe, wojskowe i administracyjne.
Nie deklaruję formalnego wykształcenia akademickiego ani statusu historyka.
Moją propozycją jest perspektywa świadka i uczestnika, która może
uzupełniać analizę wykładowcy, a nie zastępować źródła naukowe.
Spotkanie można dostosować do tematyki zajęć, poziomu grupy
i czasu przeznaczonego na dyskusję.

\end{minipage}

\vspace{0.14em}
{\color{gold}\hrule height 1pt}
\vspace{0.05em}

\begin{center}
  {\fontsize{7.1}{8.1}\selectfont\color{muted}
    Józef Szwejk
    \textbullet{} Praga
    \textbullet{} Perspektywa świadka
    \textbullet{} Opowieść i dyskusja}
\end{center}

\end{document}